\section{近期研究（2026 年 8--11 月）}
\paragraph{DynActiveGS：分别衡量结构不确定性和动态干扰}
2026.08~\cite{2608_01178}。DynActiveGS 从动态环境的 RGB-D 观测估计像素不确定性，控制观测在高斯重建中的权重，再分别用结构和运动不确定性在 Voronoi 图上选视点。规划既要补足静态几何的覆盖，也要考虑动态物体可能造成的遮挡。

论文在 Social-MP3D 与动态增强的 Gibson 模拟场景检查重建与保留视图，保留视点不出现在各方法的训练轨迹里。结果对应模拟动态环境，传感器噪声和未建模动态仍需实机验证。

\paragraph{TRACE：用连续轨迹覆盖高斯地图的信息分布}
2026.08~\cite{2608_02304}。TRACE 从在线 2D 高斯面片地图聚合体素信息与原语不确定性，再优化可执行的连续遍历轨迹。视场已经覆盖的信息随轨迹规划逐步扣减，以减少对同一位置的重复访问。

在八个 Replica 场景中相对 ActiveGS 报告平均 PSNR 改善 1.5 dB。该值适用于所比较控制规划和重建协议，信息场质量仍取决于在线地图与不确定性估计。

\paragraph{CasDeblurGS：让两幅模糊图像经二维再到三维互相约束}
2026.08~\cite{2608_10345}。相机运动使图像模糊，而且从两幅图恢复视图时，直接跨图拷贝容易将遮挡和误匹配带入结果。CasDeblurGS 先用双向光流和一致性掩码得到局部二维指引，再由冻结的无位姿高斯骨干建立临时三维场景，以重渲染输入视角提供更密集的全局指引。修复后的两幅图交给同一骨干生成可渲染高斯。

作者在 Deblur-NeRF 固定 85 个两输入一目标的三元组，真实与合成子集相对最强对照的 PSNR 分别提高 1.19 和 2.11 dB。方法在测试时不需要输入位姿或逐场优化，但需要已知内参，研究设置为静态场景。极强模糊和大遮挡会降低跨视图引导的可靠性。

\paragraph{GaussVid：用带相机条件的视频扩散修复稀疏重建伪影}
2026.08~\cite{2608_21849}。GaussVid 通过受控高斯优化构造带伪影的成对视频数据，再用相机相关条件引导视频扩散修复视图。学习目标是让已有 3DGS 的跨视角渲染更一致，生成修复的图像细节仍可能缺少真实三维观测。

论文用 DL3DV 独立场景的 500 个片段评价，并用 Zip-NeRF 大型场景测试跨数据和重建方法迁移。应把画质修复、跨视角一致性与真实三维几何分别评价。

\paragraph{FujinSplat：在 RAW 域校正烟雾颜色后再拟合高斯}
2026.09~\cite{2609_06017}。FujinSplat 在 RAW 域学习紧凑可逆的颜色调整，控制器先在可给出精确标签的模拟 RAW 数据上训练，再对真实视图预测校正。高斯优化用有界的零均值颜色残差缓解不同视图间剩余的差异。

论文在 RealX3D 的六个烟雾场景对比，ConvIR 的平均 LPIPS 为 0.493，优于本方法的 0.541，而 PSNR 低 2.57 dB。全局颜色变换无法补回传感器已经丢失的纹理，也不代表可以直接迁移到任何相机的 RAW 流程。

\paragraph{TV-SGS：让三维原语以张量投票交换几何线索}
2026.09~\cite{2609_07734}。少视图下逐原语的光度约束不足以决定真实表面。TV-SGS 用张量投票在高斯间传播局部结构线索，再将得到的三维几何监督写入损失，优化时仍保持原有渲染任务。

DTU 的十五场景比较区分三视图的大重叠和小重叠输入。补充实验的较密输入下，带有多视图几何监督的 PGSR 骨干并未因投票项明显提高。因此应按照输入重叠和骨干原有几何监督理解收益。

\paragraph{Bi-FlowGS：让补全视频与高斯几何互相修正}
2026.09~\cite{2609_17039}。稀疏视图的扩散补全如果只提供伪颜色，场景可能用错误几何解释生成图。Bi-FlowGS 先从修复视频估计光流，令高斯渲染导出的运动对应与该流对齐；更新后的高斯流和几何再回馈视频修复，构成双向细化。生成数据提供跨视图线索，但其内容不自动成为实测结构。

在 Mip-NeRF 360 的 3、6、9 输入视图条件，原文相对各配置最强对照报告 PSNR 提升 0.43、0.34 和 0.38 dB。该基准没有稠密几何真值，部分几何指标来自共同的 COLMAP 轨迹。应区别渲染增益、与稀疏几何的一致性，以及没有条件确认的生成区域。

\paragraph{CADSplat：利用近似 CAD 对齐并约束少视图高斯}
2026.09~\cite{2609_18473}。CADSplat 检索一个近似 CAD 形状，用多视图轮廓对齐相机，并在高斯外观拟合时联合细化注册与非刚性形变。CAD 为不可见部分提供形状约束，图像仍决定可观测的外观和局部修正。

作者在 SCO-CAD、NeRS MVMC 和 NeRS Misc 上评价。主表的 SCO-CAD 保留视角结果含测试时位姿细化，部分基线未使用该细化。方法依赖合理的初始相对标定及近似可用的形状先验。

\paragraph{用适配后的前馈几何给稀疏水下视图补足匹配证据}
2026.09~\cite{2609_18737}。水下散射和吸收会同时污染颜色及几何伪标签。该方法保留清晰图的几何教师，用物理退化后的图训练 VGGT 学生，只给视觉聚合编码器装入 LoRA 并冻结预测头。适配后的 VGGT 从未知位姿输入预测稠密点，其高斯代理渲染相邻伪视图与深度，再让真实视图和伪视图共同参与 COLMAP，所得轨迹与深度用于初始化 RUSplatting。

原文在 SeaThru-NeRF 和 Submerged3D 各选四个场景，每八帧留出一帧测试。SeaThru-NeRF 上，整体流程将 COLMAP 初始化的 RUSplatting 平均 PSNR 从 24.37 提高到 27.11 dB。IUI3-RedSea 的强散射仍会破坏 SIFT 匹配，少视角下也有残余预测噪声，因此不能将两组数据上的结果外推到所有浑浊水域。

\paragraph{D3GS：用米制深度、跨视图特征和扩散辅助少视图优化}
2026.09~\cite{2609_22941}。D3GS 首先把来自 MapAnything 与束调整的稀疏米制深度补全后初始化高斯，再用深度、DINO 特征与扩散修复为优化提供互补监督。几何先验、跨视图对齐和生成外观分别约束不同误差来源。

论文报告少视图设定下的几何和画质比较。若初始米制深度失准，后续初始化可能继承误差；扩散生成补充的图像外观也需要通过跨视图观测判断是否对应真实几何。

\paragraph{GAPS：按信息价值和生成可靠度筛选伪视图}
2026.09~\cite{2609_23436}。GAPS 在少视图高斯的交替优化中，用深度、外观和场景适配约束生成器，再按候选视角的重建信息和生成可靠度安排补充，从保守内插逐步过渡到外推。通过检查的伪视图按像素不确定性进入训练，高密度区域另加高斯 dropout。

论文在 LLFF 和 Mip-NeRF 360 使用固定保留比例评价。LLFF 六视图相对原始 3DGS 平均 PSNR 高 0.89 dB。消融中移除视角选择会让 Mip-NeRF 360 的 LPIPS 劣于无伪视图的基线，生成先验的适用范围因此与筛选质量和场景覆盖有关。

\paragraph{LiTe-GS：用随机候选集减少主动选视的信息计算}
2026.09~\cite{2609_30393}。主动重建每次选下一视角时，如果对所有候选都计算 Fisher 信息，选视本身会成为开销。LiTe-GS 保留 POP-GS 的信息目标，在每轮贪心选择时只评估随机抽出的候选子集。近似参数控制样本量，选中的帧仍进入原有的逐场景高斯优化。

作者在 Blender 和 Mip-NeRF 360 上比较图像质量和 Fisher 评分调用次数。调用数下降而重建表现接近穷举选视的条件随子集大小而变。该效率结论度量选视阶段的调用数，不能直接视为端到端训练时间的同等降幅。

\paragraph{用双向轨迹可靠度调节渐进式无 COLMAP 重建}
2026.09~\cite{2609_30865}。该方法用 SE(3) 上的短时运动连续性提供相机位姿初值，再用正反向跟踪一致性衡量置信度，为渐进式重建中的传播和回溯调整提供控制信号。置信度用于筛选轨迹证据并减弱错误早期位姿向后传播。

论文把相机误差和保留视图画质分别列出，前者用 ATE 与相对位姿误差，后者用 PSNR、SSIM 等。双向一致性是内部可靠性线索，不能替代在急剧视点变化和跟踪失败时的外部位姿参考。

\paragraph{ClearGS：按可靠度保留手持视频的几何覆盖}
2026.09~\cite{2609_31509}。手持视频的模糊分布不均，直接删除低质量帧可能丢失某条轨迹上的全部观察。ClearGS 将视觉可靠度、退化风险和几何效用用于分配不同训练权重，保留弱覆盖区域的监督。修复以当前渲染作为位姿对齐的结构约束，以同一帧的原观测作为待恢复证据，最后汇总沿轨迹接受的修复。

论文在 GS2E 的三个模糊等级及 GSOTM 的模糊、滚动快门、两者结合和位姿噪声条件上测试。训练与修复过程没有配对清晰图，清晰目标只用于评价。感知质量改善不能证明修复得到的纹理都来自场景，作者也将修复证据的几何核对列为后续工作。

\paragraph{OC-GS：用固定轴运动模型校正不规则转台角度}
2026.09~\cite{2609_31572}。OC-GS 从不规则间隔的转台图像获得初始角度，在共享固定转轴的运动约束下联合细化角度和高斯几何。它利用的是可控的采集结构，所以能减少对一般自由运动相机估计的依赖。

在相同训练器和初始化的消融里，十二个不规则视图下细化角度比固定初始角度高 9.08 dB 前景 PSNR，六图时高 4.22 dB。该实验对应固定转轴扰动，方法仍依赖 VGGT 初始化和较小的转轴噪声。

\paragraph{AGILE-GS：用虚拟锚点和短名单减少主动视角筛选成本}
2026.09~\cite{2609_34176}。AGILE-GS 先在连续位姿空间优化虚拟锚点，以代表地图当前的信息方向，再用低冗余的候选短名单替代全候选渲染评估。基础模式避免逐候选 Fisher 重排，增强模式在短名单中加入重排。

作者在 Mip-NeRF 360、Tanks and Temples、Deep Blending 分开报告平均选点时间与重建画质，并另给具身采集结果。选点阶段的节省不等于相同比例的完整采集缩短，真实运动时间还由采集协议与可执行路径决定。

\paragraph{UGOD：用原语不确定性调节透明度和训练 dropout}
2026.09~\cite{2609_39089}。UGOD 从高斯属性与观察方向估计视图相关不确定性。它在透明度合成前连续抑制高不确定贡献，并以梯度分离的软 dropout 正则训练中的不可靠原语。前者作用于渲染，后者作用于优化。

论文在 LLFF 三视图及 Mip-NeRF 360 的八和二十四视图设定评价。LLFF 使用 COLMAP 重建的三个训练视图，不能把结果直接推断给未知相机参数或其他稀疏度。训练细节和画质结果应按协议对应。

\paragraph{MoonGS：从双视图和语义先验重建低纹理月面}
2026.10~\cite{2610_07110}。MoonGS 从一对带相机参数和语义图像的月面视图直接预测高斯。方法利用视觉基础模型的深度特征确定基元位置，在非重叠或缺少纹理的区域以语义先验改善深度估计，并用深度排序约束稳定背景。启发式视图重采样用于补足稀疏观测。

作者在 LuSNAR 和合成 MoonBlender 上以重训的前馈新视图模型为基线，LuSNAR 的 PSNR 由所比基线的最高 18.89 dB 提高到 23.80 dB。真实嫦娥图像只用于定性比较。该增益不能直接解读为真实月面的米制几何误差下降，两幅输入也限制了大面积遮挡区的证据。

\paragraph{位姿图替代束调整的失效分析}
2026.10~\cite{2610_14090}。作者检验密集闭环下用仅包含位姿的图优化替代全局束调整。实验中的子块仍需做局部束调整，跨块再用位姿约束合并。纯 SE(3) 边不能表达各子块尺度差异，成对边也不能恢复同一个三维点在多视图间形成的联合约束。

在两个详细展示的自采无人机场景及公开 GauU-Scene V2 中，所测实现未达到调优 COLMAP 精度。消费级 GPS、依序合并和手工权重也可能造成差距。由于没有对这些选择分别做消融，论文的负结论不能简单扩大到所有分布式估计或保留稀疏结构的混合方法。

\paragraph{以少样本训练控制动态稀疏视图高斯过拟合}
2026.10~\cite{2610_25887}。作者把动态稀疏视图重建的过拟合拆为监督不足、相机偏差和模型容量三个问题。光流用于合成额外监督视图，重投影误差用于联合调整相机与场景。训练端选择性丢弃静态高斯，并以早停减少对少量观测的共同适配。

定量表来自 SIGGRAPH Asia 2025 稀疏视角竞赛的两个验证序列，每个用八个训练视角与八个新视角评价，只计前景区域。伪视角生成假定物体运动适度，特别快的人体动作可能破坏该条件。

\paragraph{CrossFit-GS：用目标帧未参与训练的子图精调位姿}
2026.10~\cite{2610_28190}。逐帧推进的无 COLMAP 高斯重建可能把某帧的位姿误差吸收到高斯外观中，再用同一地图确认错误位姿。CrossFit-GS 暂缓把候选位姿写入全局地图，先以未使用目标帧且参数冻结的子图精化位姿，通过另留出的像素决定是否采用结果。

作者在 Tanks and Temples 与 CO3D-V2 同时评价轨迹和新视图。消融表明，对比只在原图上降低光度损失，隔离目标观测更有助于避免误差回流。交叉拟合仍需有足够重叠以建立可用参考，并要保守处理覆盖不足的更新。

\paragraph{DomeGS：用带视差的高斯天空穹顶完善 VR 远景}
2026.10~\cite{2610_37172}。DomeGS 从已有室外高斯场抽取全景，以考虑全周一致性的生成外扩补全背景，再把全景转成几何适配的多视角监督。背景最终作为天空穹顶高斯与原场景共同渲染，头部移动时可以产生经典无限远天空盒缺少的视差。

作者用 36 人的被试内 VR 研究比较天空盒与高斯穹顶，主要任务引导参与者观察地平线和上半球。结果支持测试场景下的背景一致性与临场感，补出的远景不是实际测量几何。

\paragraph{谱系感知稀疏高斯训练：阶段式随机屏蔽与均衡增密}
2026.11~\cite{2611_17752}。克隆及分裂可能产生协同拟合同一训练视角的一族高斯。该工作以增密谱系记录原语关系，早期对整个谱系共享随机屏蔽，随后改为独立屏蔽；同时在不改变克隆与分裂总配额的前提下，调整各谱系获得的增密机会。

论文在 LLFF、Mip-NeRF 360 与 Blender 中比较，并测试把两个模块接入不同主干的效果。作者报告收益因场景、指标和主干而异，阶段式屏蔽是主要贡献，均衡分配提供额外改进。

\paragraph{QuantileGS：为冻结高斯渲染器校准新视角误差上界}
2026.11~\cite{2611_21253}。QuantileGS 固定已有三维高斯的 RGB 渲染，再给原语拟合有序误差分位属性。留出的校准视角估计归一化校准量，从而为新视角给出随位置变化的 RGB 平均绝对误差上界。查询需要相机参数，不需要目标图像。

作者在 LF、LLFF 和 Mip-NeRF 360 比较误差排序与覆盖，并展示主动选视、裁剪和 SLAM 的经验收益。有限样本覆盖是在校准与测试误差分数可交换时对固定渲染器成立的边际保证，不能扩展为每个像素、每个视图同时有保证。
